\documentclass[10pt,a4paper]{amsart}
\usepackage[margin=19mm]{geometry}
\usepackage{amsmath,amssymb,mathtools}
\usepackage[scr=rsfs,cal=euler]{mathalfa}
\usepackage{xfrac}
\usepackage[unicode,hidelinks,bookmarks=false]{hyperref}
\newcommand{\ie}{i.e.\ }
\newcommand{\cf}{cf.\ }
\newcommand{\resp}{respectively\ }
\newcommand{\Subsection}{\subsection}
\providecommand{\nobreakdash}{\nobreak\hbox{-}}
\setlength{\emergencystretch}{2em}
\multlinegap0pt
\usepackage{bm}
\usepackage{bbm}
\usepackage{dsfont}
\usepackage{xspace}
\usepackage{cancel}
\usepackage{mathtools}
\usepackage{amsthm}
\makeatletter
\@ifundefined{thm}{\newtheorem{thm}{Thm}}{}
\@ifundefined{lem}{\newtheorem{lem}[thm]{Lemma}}{}
\makeatother
\DeclareMathOperator{\id}{id}
\def\cal{\mathcal}
\newcommand{\Prob}{\mathrm{Prob}}
\title[Existentially closed measure-preserving actions of approxima]{Existentially closed measure-preserving actions of approximately treeable groups}
\author{Isaac Goldbring, Brandon Seward,, Robin Tucker-Drob}
\date{Source: arXiv:2507.03195v1 (CC BY 4.0)}
\begin{document}
\maketitle

\noindent\emph{Source and licence.} Existentially closed measure-preserving actions of approximately treeable groups. Version arXiv:2507.03195v1 (CC BY 4.0), distributed under the Creative Commons Attribution 4.0 International licence, as recorded in the arXiv licence metadata for this exact version and shown on the abstract page https://arxiv.org/abs/2507.03195v1. Full rights notice: licenses/arxiv-2507.03195-CC-BY-4.0.txt.\par\smallskip

\noindent\textit{Editorial scope.} Counted unit: the Lem whose source label is \texttt{lem:fmap} in the article (source0.tex lines 1997--2004, source label \texttt{lem:fmap}), delivered together with the article's own complete original proof of it (source0.tex lines 2006--2097, one \texttt{proof} environment of 92 lines). No mathematics is written, repaired or re-derived: statement and proof are the article's own text line for line. Required context retained uncounted: lem:comp (lines 1892--1894); lem:mfold (lines 1911--1915); lem:pairmap (lines 1928--1936); item:pmap1 (lines 1930--1935); item:pmap2 (lines 1930--1935); item:pmap3 (lines 1930--1935). Every definition, display and label of those lines is retained unchanged. Omitted from this package and not redistributed: 1--1891, 1895--1910, 1916--1927, 1936--1996, 2005--2005, 2098--2229. Nothing inside a retained excerpt is omitted or altered: every retained body is delivered exactly as the article's lines are written, so no omission receipt of any kind accompanies this package. Citations: the retained excerpts carry no citation command at all, so no bibliography entry of the article is transcribed and no reference is redistributed. Reference adaptation: none was needed and none was made; every reference inside the delivered unit resolves inside it, so no unresolved reference or citation is produced. Printed slips kept literal: no printed word, symbol, hypothesis or number of the counted statement or of its proof was altered. Layout and class: the article's own class is \texttt{amsart} with packages including tgpagella, euler, fontenc, amsmath, amssymb, hyperref, babel, mathrsfs, eucal, xy, tikz, enumitem, mathtools; the wrapper is the lane's pinned \texttt{amsart} editorial wrapper at 10pt on A4, which restates the theorem-like environments the retained text needs and adds the article's own macro definitions that the retained text needs (\texttt{\textbackslash Prob}, \texttt{\textbackslash cal}, \texttt{\textbackslash id}), with extra wrapper packages (bm, bbm, dsfont, xspace, cancel, mathtools). The delivered numbering is the unit's own. Assets: no retained span contains a figure environment, a graphics-inclusion command, a PDF-page inclusion, a table or a TikZ picture; no image file is copied into this package, the package declares no asset dependency and no image file is redistributed with it. Licence: version arXiv:2507.03195v1 of this article is distributed under the Creative Commons Attribution 4.0 International licence, as recorded in the arXiv licence metadata for this exact version and shown on the abstract page https://arxiv.org/abs/2507.03195v1. A full rights notice is delivered at licenses/arxiv-2507.03195-CC-BY-4.0.txt. Characters: every retained excerpt is pure ASCII, so the delivered engine, XeLaTeX with its default Latin Modern text font, reports no missing character.
\begin{lem} \label{lem:comp}
	The composition function $((h_1, \omega_1),(h_0,\omega_0)) \in M_2 \mapsto (h_1 \circ h_0, \omega_0) \in M$ is continuous.
\end{lem}

\begin{lem} \label{lem:mfold}
	The map
	\[((h_\gamma)_{\gamma \in C}, \omega) \in M^{(C)} \mapsto \left(\prod_{\gamma \in C} h_\gamma \right)_* \omega \in \Prob((p^\Gamma)^C)\]
	is continuous.
\end{lem}

\begin{lem} \label{lem:pairmap}
	There exists a function $\psi$ assigning to each pair $(\kappa_1, \kappa_0) \in P^*$ a Borel measurable function $\psi(\kappa_1, \kappa_0) : p^\Gamma \rightarrow p^\Gamma$ such that
	\begin{enumerate}
		\item \label{item:pmap1} $\kappa_1 = \psi(\kappa_1, \kappa_0)_* \kappa_0$;
		\item \label{item:pmap2} $\psi(\kappa_0, \kappa_1) \circ \psi(\kappa_1, \kappa_0)$ is equal to the identity $\kappa_0$-almost-everywhere;
		\item \label{item:pmap3} the map $(\kappa_1, \kappa_0) \in P^* \mapsto (\psi(\kappa_1, \kappa_0), \kappa_0) \in M$ is continuous; and
		\item \label{item:pmap4} $(\psi(\kappa_1, \kappa_0), \kappa_0) \rightarrow (\id, \kappa)$ as $(\kappa_1, \kappa_0) \rightarrow (\kappa, \kappa)$.
	\end{enumerate}
\end{lem}

\begin{lem} \label{lem:fmap}
	There is a map $\theta : P^{(\Gamma)} \times \cal F(\Gamma) \rightarrow \Prob(p^\Gamma)$ such that:
	\begin{enumerate}
		\item[(a)] $\theta(\omega, F)$ is a Borel function of $(\omega, F)$ and is a continuous function of $\omega$;
		\item[(b)] $\theta(\omega, F) = \phi(\omega(e), E_F)$ whenever $\omega \in P^{(\Gamma)}$ satisfies $\gamma^s_* \omega(\gamma) = \omega(e)$ for all $\gamma \in \Gamma$;
		\item[(c)] $\gamma^s_* \theta(\omega, F) = \theta(\gamma^s \cdot \omega, \gamma^d \cdot F)$ for all $\gamma \in \Gamma$, $\omega \in P^{(\Gamma)}$ and $F \in \cal F(\Gamma)$.
	\end{enumerate}
\end{lem}

\begin{proof}
	Let $\psi$ be the function described in Lemma \ref{lem:pairmap}. For $\delta, \gamma \in \Gamma$ and $\omega \in P^{(\Gamma)}$ define Borel functions $g_{\delta, \gamma}^\omega, \bar g_{\delta, \gamma}^\omega : p^\Gamma \rightarrow p^\Gamma$ by
	\[g_{\delta, \gamma}^\omega = \psi(\omega(\delta), (\delta^{-1} \gamma)^s_* \omega(\gamma)) \circ (\delta^{-1} \gamma)^s\]
	and
	\[\bar g_{\delta, \gamma}^\omega = (\delta^{-1} \gamma)^s \circ \psi((\gamma^{-1} \delta)^s_* \omega(\delta), \omega(\gamma)).\]
	The definition of $P^{(\Gamma)}$ ensures that $\psi$ is defined in both of the expressions above. With these definitions, $g_{\delta, \gamma}^\omega$ first performs the shift on $p^\Gamma$ by $\delta^{-1} \gamma$ and then, according to $\psi$, applies a map from $p^\Gamma$ to itself that pushes $(\delta \gamma^{-1})_* \omega(\gamma)$ forward to $\omega(\delta)$. Similarly, $\bar g_{\delta, \gamma}^\omega$ first applies a map from $p^\Gamma$ to itself that pushes $\omega(\gamma)$ forward to $(\gamma^{-1} \delta)^s_* \omega(\delta)$ and then performs the shift on $p^\Gamma$ by $\delta^{-1} \gamma$.
	
	We observe that the following statements hold for all $\delta, \gamma \in \Gamma$:
	\begin{enumerate}
		\item \label{proofitem:fmap1} Each map $g_{\delta,\gamma}^\omega$ and $\bar g_{\delta, \gamma}^\omega$ pushes $\omega(\gamma)$ forward to $\omega(\delta)$. This is clear from the definitions together with clause (\ref{item:pmap1}) of Lemma \ref{lem:pairmap}.
		\item \label{proofitem:fmap2} \indent Both compositions $\bar g_{\gamma, \delta}^\omega \circ g_{\delta, \gamma}^\omega$ and $g_{\gamma, \delta}^\omega \circ \bar g_{\delta, \gamma}^\omega$ are equal to the identity $\omega(\gamma)$-almost-everywhere. This follows from the definitions and clause (\ref{item:pmap2}) of Lemma \ref{lem:pairmap}.
		\item \label{proofitem:fmap3} The pairs $(g_{\delta, \gamma}^\omega, \omega(\gamma)), (\bar g_{\delta, \gamma}^\omega, \omega(\gamma)) \in M$ are continuous functions of $\omega$. Indeed each pair can be expressed as the result of a composition $M_2 \rightarrow M$:
		\[(g_{\delta, \gamma}^\omega, \omega(\gamma)) = (\psi(\omega(\delta), (\delta^{-1} \gamma)^s_* \omega(\gamma)), (\delta^{-1} \gamma)^s_* \omega(\gamma)) \cdot ((\delta^{-1} \gamma)^s, \omega(\gamma))\]
		\[(\bar g_{\delta, \gamma}^\omega, \omega(\gamma)) = ((\delta^{-1} \gamma)^s, (\gamma^{-1} \delta)^s_* \omega(\delta)) \cdot (\psi((\gamma^{-1} \delta)^s_* \omega(\delta), \omega(\gamma)), \omega(\gamma)).\]
		Continuity is a consequence of Lemma \ref{lem:comp}, Lemma \ref{lem:pairmap}.(\ref{item:pmap3}), and the continuity of the push-forward maps $\gamma^s_* : \Prob(p^\Gamma) \rightarrow \Prob(p^\Gamma)$ for $\gamma \in \Gamma$.
		\item \label{proofitem:fmap4} When $\delta^s_* \omega(\delta) = \gamma^s_* \omega(\gamma)$ both $g_{\delta, \gamma}^\omega$ and $\bar g_{\delta, \gamma}^\omega$ are equal to $(\delta^{-1} \gamma)^s$.
		\item \label{proofitem:fmapa} $g_{\beta \delta, \beta \gamma}^{\beta^s \cdot \omega} = g_{\delta, \gamma}^\omega$ and $\bar g_{\beta \delta, \beta \gamma}^{\beta^s \cdot \omega} = \bar g_{\delta, \gamma}^\omega$ for all $\beta \in \Gamma$. This is immediate from the definitions together with the facts that $(\beta^s \cdot \omega)(\beta \delta) = \omega(\delta)$, $(\beta^s \cdot \omega)(\beta \gamma) = \omega(\gamma)$, and $(\beta \delta)^{-1} (\beta \gamma) = \delta^{-1} \gamma$.
	\end{enumerate}
	
	Now for $\omega \in P^{(\Gamma)}$ and $F \in \cal F(\Gamma)$ we define a function $\sigma_F^\omega$ assigning to each pair $(\delta, \gamma) \in E_F$ a Borel measurable function $\sigma_F^\omega(\delta, \gamma) : p^\Gamma \rightarrow p^\Gamma$ as follows. First, when $\delta = \gamma$ we set $\sigma_F^\omega(\delta, \gamma)$ equal to the identity function. Next, if $(\delta, \gamma) \in F \cup \bar F$ then we set
	\[\sigma_F^\omega(\delta, \gamma) = \begin{cases}
		g_{\delta, \gamma}^\omega & \text{if } (\delta, \gamma) \in F\\
		\bar g_{\delta, \gamma}^\omega & \text{if } (\delta, \gamma) \in \bar F.
	\end{cases}\]
	In general for $(\delta, \gamma) \in E_F$, consider the unique sequence of vertices $\gamma_0, \gamma_1, \ldots, \gamma_n$ forming a path in $F \cup \bar F$ from $\gamma$ to $\delta$, specifically $\gamma_0 = \gamma$, $\gamma_n = \delta$, and $(\gamma_{i+1}, \gamma_i) \in F \cup \bar F$ for all $0 \leq i < n$, and set
	\[\sigma_F^\omega(\delta, \gamma) = \sigma_F^\omega(\gamma_n, \gamma_{n-1}) \circ \cdots \circ \sigma_F^\omega(\gamma_1, \gamma_0).\]
	
	The following properties hold for all $\omega \in P^{(\Gamma)}$, $F \in \cal F(\Gamma)$, and $(\delta, \gamma) \in E_F$.
	\begin{enumerate}\setcounter{enumi}{5}
		\item \label{proofitem:fmap5} $\sigma_F^\omega(\delta, \gamma)$ pushes $\omega(\gamma)$ forward to $\omega(\delta)$. This is an immediate consequence of (\ref{proofitem:fmap1}) above.
		\item \label{proofitem:fmap6} If $\beta$ lies on the path from $\gamma$ to $\delta$ then $\sigma_F^\omega(\delta, \gamma) = \sigma_F^\omega(\delta, \beta) \circ \sigma_F^\omega(\beta, \gamma)$. This is immediate from the definition.
		\item \label{proofitem:fmap7} $\sigma_F^\omega(\gamma, \delta) \circ \sigma_F^\omega(\delta, \gamma)$ is equal to the identity $\omega(\gamma)$-almost-everywhere. When $(\delta, \gamma) \in F \cup \bar F$ this fact follows immediately from (\ref{proofitem:fmap2}) above. The general case follows by induction on the length of the path from $\gamma$ to $\delta$. Specifically, if $\beta$ lies on the path from $\gamma$ to $\delta$ then by (\ref{proofitem:fmap6})
		\[\sigma_F^\omega(\gamma, \delta) \circ \sigma_F^\omega(\delta, \gamma) = \sigma_F^\omega(\gamma, \beta) \circ \sigma_F^\omega(\beta, \delta) \circ \sigma_F^\omega(\delta, \beta) \circ \sigma_F^\omega(\beta, \gamma).\]
		By induction we can assume $\sigma_F^\omega(\beta, \delta) \circ \sigma_F^\omega(\delta, \beta)$ is equal to the identity $\omega(\beta)$-almost-everywhere. Since $\sigma_F^\omega(\beta, \gamma)$ pushes $\omega(\gamma)$ forward to $\omega(\beta)$ by (\ref{proofitem:fmap5}), it follows that $\omega(\gamma)$-almost-everywhere $\sigma_F^\omega(\gamma, \delta) \circ \sigma_F^\omega(\delta, \gamma)$ is equal to $\sigma_F^\omega(\gamma, \beta) \circ \sigma_F^\omega(\beta, \gamma)$, and by induction this latter function is equal to the identity $\omega(\gamma)$-almost-everywhere.
		\item \label{proofitem:fmap8} Whenever $\beta$ is in the same connected component as $\delta$ and $\gamma$, $\sigma_F^\omega(\delta, \gamma)$ is $\omega(\gamma)$-almost-everywhere equal to $\sigma_F^\omega(\delta, \beta) \circ \sigma_F^\omega(\beta, \gamma)$. To see this, let $\beta'$ be the unique point belonging to all three of the paths between the points $\delta$, $\gamma$, and $\beta$. Two applications of (\ref{proofitem:fmap6}) yield the two equations:
		\begin{align*}
			\sigma_F^\omega(\delta, \beta) \circ \sigma_F^\omega(\beta, \gamma) & = \sigma_F^\omega(\delta, \beta') \circ \sigma_F^\omega(\beta', \beta) \circ \sigma_F^\omega(\beta, \beta') \circ \sigma_F^\omega(\beta', \gamma)\\
			\sigma_F^\omega(\delta, \gamma) & = \sigma_F^\omega(\delta, \beta') \circ \sigma_F^\omega(\beta', \gamma).
		\end{align*}
		The right-sides of these two equations are $\omega(\gamma)$-almost-everywhere equal since $\sigma_F^\omega(\beta', \gamma)$ pushes $\omega(\gamma)$ forward to $\omega(\beta')$ by (\ref{proofitem:fmap5}) and $\sigma_F^\omega(\beta', \beta) \circ \sigma_F^\omega(\beta, \beta')$ is $\omega(\beta')$-almost-everywhere equal to the identity by (\ref{proofitem:fmap7}). This establishes the claim.
		\item \label{proofitem:fmap10} If $\gamma^s_* \omega(\gamma) = \omega(e)$ for all $\gamma \in \Gamma$ then $\sigma_F^\omega(\delta, \gamma) = (\delta^{-1} \gamma)^s$ for all $(\delta, \gamma) \in E_F$. This follows from (\ref{proofitem:fmap4}).
		\item \label{proofitem:fmapb} $\sigma_{\beta^d \cdot F}^{\beta^s \cdot \omega}(\beta \delta, \beta \gamma) = \sigma_F^\omega(\delta, \gamma)$ for all $\beta \in \Gamma$. This is a consequence of (\ref{proofitem:fmapa}) together with the fact that if $\gamma_0, \ldots, \gamma_n$ is the sequence of vertices forming a path in $F$ from $\gamma$ to $\delta$ then $\beta \gamma_0, \ldots, \beta \gamma_n$ is the sequence of vertices for the path from $\beta \gamma$ to $\beta \delta$ in $\beta^d \cdot F$.
	\end{enumerate}
	Additionally, we observe the following facts:
	\begin{enumerate}\setcounter{enumi}{10}
		\item \label{proofitem:fmapc} The function $(\omega, F) \in P^{(\Gamma)} \times U_{\delta, \gamma} \mapsto (\sigma^\omega_F(\delta, \gamma), \omega(\gamma)) \in M$ is continuous, where $U_{\delta, \gamma}$ is the open set $\{F \in \cal F(\Gamma) : (\delta, \gamma) \in E_F\}$. Specifically, we can partition $U_{\delta, \gamma}$ into a collection $\cal U$ of clopen sets so that for every $U \in \cal U$ the path from $\delta$ to $\gamma$ in $F$ is constant over all $F \in U$. It's enough to observe that for each $U \in \cal U$ the map $(\omega, F) \mapsto (\sigma^\omega_F(\delta, \gamma), \omega(\gamma))$ is continuous on $P^{(\Gamma)} \times U$. Say $\gamma = \gamma_0, \gamma_1, \ldots, \gamma_n = \delta$ is the path from $\gamma$ to $\delta$ in $F$ for all $F \in U$. Then from the definition of $\sigma^\omega_F$ and by (\ref{proofitem:fmap5}) we see that $(\sigma^\omega_F(\delta, \gamma), \omega(\gamma)) \in M$ is the composition of the elements $(\sigma^\omega_F(\gamma_{i+1}, \gamma_i), \omega(\gamma_i)) \in M$ when $F \in U$. It then follows from (\ref{proofitem:fmap3}) and Lemma \ref{lem:comp} that $(\sigma^\omega_F(\delta, \gamma), \omega(\gamma))$ is a continuous function of $(\omega, F) \in P^{(\Gamma)} \times U$.
		\item \label{proofitem:fmap9} If $C \subseteq \Gamma$ is nonempty, $\gamma \in \Gamma$, and $C \cup \{\gamma\}$ is contained in a single $E_F$-class for all $F$ in a set $H \subseteq \cal F(\Gamma)$, then the function $(\omega, F) \in P^{(\Gamma)} \times H \mapsto ((\sigma_F^\omega(\delta, \gamma))_{\delta \in C}, \omega(\gamma)) \in M^{(C)}$ is continuous. Indeed, for every $\delta \in C$ the map $(\omega, F) \in P^{(\Gamma)} \times H \mapsto (\sigma_F^\omega(\delta, \gamma), \omega(\gamma)) \in M$ is continuous by (\ref{proofitem:fmapc}) since $H \subseteq U_{\delta, \gamma}$.
	\end{enumerate}
	
	Let $C \in \Gamma / E_F$ be a connected component and fix any $\gamma_C \in C$. From $\omega$ we define a Borel probability measure on $(p^\Gamma)^C$ by the formula
	\[\left( \prod_{\delta \in C} \sigma_F^\omega(\delta, \gamma_C) \right)_* \omega(\gamma_C).\]
	Notice that this measure depends continuously on $\omega$ by (\ref{proofitem:fmap9}) and Lemma \ref{lem:mfold}. Additionally, the measure we obtain does not depend upon the choice of $\gamma_C \in C$ since for any other element $\gamma_C' \in C$ properties (\ref{proofitem:fmap5}) and (\ref{proofitem:fmap8}) imply
	\begin{align*}
		\left(\prod_{\delta \in C} \sigma_F^\omega(\delta, \gamma_C) \right)_* \omega(\gamma_C) & = \left(\prod_{\delta \in C} \sigma_F^\omega(\delta, \gamma_C) \right)_* \sigma_F^\omega(\gamma_C, \gamma_C')_* \omega(\gamma_C')\\
		& = \left( \prod_{\delta \in C} \sigma_F^\omega(\delta, \gamma_C) \circ \sigma_F^\omega(\gamma_C, \gamma_C') \right)_* \omega(\gamma_C')\\
		& = \left( \prod_{\delta \in C} \sigma_F^\omega(\delta, \gamma_C') \right)_* \omega(\gamma_C').
	\end{align*}
	
	For each $C \in \Gamma / E_F$ pick some $\gamma_C \in C$. Also let $f : (p^\Gamma)^\Gamma \rightarrow p^\Gamma$ be the flattening map given by the formula $f(z)(\gamma) = z(\gamma)(e)$. We define the Borel probability measure $\theta(\omega, F) \in \Prob(p^\Gamma)$ by
	\[\theta(\omega, F) = f_* \prod_{C \in \Gamma / E_F} \left( \prod_{\delta \in C} \sigma_F^\omega(\delta, \gamma_C) \right)_* \omega(\gamma_C).\]
	Formally the expression above to the right of $f_*$ is a product measure on a product space indexed by the set $\Gamma / E_F$, however since the measure associated to the coordinate $C \in \Gamma / E_F$ is a measure on $(p^\Gamma)^C$, we identify the expression to the right of $f_*$ as a Borel probability measure on $(p^\Gamma)^\Gamma$ (and therefore the application of $f_*$ is defined).
	
	We first check (b). Denote by $\pi_C : p^\Gamma \rightarrow p^C$ the projection map for $C \subseteq \Gamma$, and recall that
	\[\phi(\omega(e), E_F) = \prod_{C \in \Gamma / E_F} (\pi_C)_* \omega(e).\]
	Let us also write $f$ for the map from $(p^\Gamma)^C$ to $p^C$ given by the same formula as before: $f(z)(\gamma) = z(\gamma)(e)$ for $\gamma \in C$. Since $\prod_{\gamma \in C} (\gamma^{-1})^s$ maps $p^\Gamma$ to $(p^\Gamma)^C$, it is easily checked that $\pi_C = f \circ \prod_{\gamma \in C} (\gamma^{-1})^s$. If $\gamma^s_* \omega(\gamma) = \omega(e)$ for all $\gamma \in \Gamma$ then using (\ref{proofitem:fmap10}) we obtain
	\begin{align*}
		\prod_{C \in \Gamma / E_F} (\pi_C)_* \omega(e) & = \prod_{C \in \Gamma / E_F} f_* \left( \prod_{\delta \in C} (\delta^{-1})^s \right)_* \omega(e)\\
		& = \prod_{C \in \Gamma / E_F} f_* \left( \prod_{\delta \in C} (\delta^{-1})^s \right)_* (\gamma_C)^s_* \omega(\gamma_C)\\
		& = \prod_{C \in \Gamma / E_F} f_* \left( \prod_{\delta \in C} (\delta^{-1} \gamma_C)^s \right)_* \omega(\gamma_C)\\
		& = f_* \prod_{C \in \Gamma / E_F} \left( \prod_{\delta \in C} \sigma^\omega_F(\delta, \gamma_C) \right)_* \omega(\gamma_C) = \theta(\omega, F),\\
	\end{align*}
	confirming (b).
	
	Next we check (a). Let $\Delta \subseteq \cal F(\Gamma) \times \cal E(\Gamma)$ be the Borel set $\Delta = \{(F, E_F) : F \in \cal F(\Gamma)\}$. Since the map $F \in \cal F(\Gamma) \mapsto (F, E_F) \in \Delta$ is Borel, it suffices to show that $\theta(F, \omega)$ is a continuous function of $(\omega, F, E_F) \in P^{(\Gamma)} \times \Delta$. In fact it is enough to show that $\theta(\omega, F)(A)$ is a continuous function of $(\omega, F, E_F)$ when $A$ is a cylinder set. Say $A = \prod_{\gamma \in \Gamma} B_\gamma$ where $\varnothing \neq B_\gamma \subseteq p$ for all $\gamma \in \Gamma$ and $B_\gamma = p$ for all $\gamma \in \Gamma \setminus D$ where $D \subseteq \Gamma$ is finite. Partition $\Delta$ into a collection $\cal V$ of clopen sets so that for each $V \in \cal V$ the restriction of $E_F$ to $D$ is constant over all $(F, E_F) \in V$. It will be enough to check that for each $V \in \cal V$ the map $(\omega, F, E_F) \mapsto \theta(\omega, F)(A)$ is continuous on $P^{(\Gamma)} \times V$. Say every $(F, E_F) \in V$ partitions $D$ into the family of sets $\cal D$. When $(F, E_F) \in V$ and $C \in \Gamma / E_F$ we have that
	\[\left(f_* \left( \prod_{\delta \in C} \sigma^\omega_F(\delta, \gamma_C) \right)_* \omega(\gamma_C) \right) \left(\prod_{\gamma \in C} B_\gamma\right)\]
	is equal to $1$ if $C \cap D = \varnothing$ and is equal to
	\[\left(f_* \left( \prod_{\delta \in C'} \sigma^\omega_F(\delta, \gamma_{C'}) \right)_* \omega(\gamma_{C'}) \right) \left(\prod_{\gamma \in C'} B_\gamma \right)\]
	when $C \cap D = C' \in \cal D$ and $\gamma_{C'}$ is any element of $C'$ (changing $\gamma_C$ to $\gamma_{C'}$ does not change the resulting value, as explained three paragraphs prior). Now notice that the map
	\[(\omega, F, E_F) \in P^{(\Gamma)} \times V \mapsto f_* \left( \prod_{\delta \in C'} \sigma^\omega_F(\delta, \gamma_{C'}) \right)_* \omega(\gamma_{C'}) \in \Prob(p^{C'})\]
	is continuous by (\ref{proofitem:fmap9}), Lemma \ref{lem:mfold}, and the fact that $f_*$ is continuous. Therefore when $(\omega, F, E_F) \in P^{(\Gamma)} \times V$ we have that
	\[\theta(\omega, F)(A) = \prod_{C' \in \cal D} \left(f_* \left( \prod_{\delta \in C'} \sigma^\omega_F(\delta, \gamma_{C'}) \right)_* \omega(\gamma_{C'}) \right) \left(\prod_{\gamma \in C'} B_\gamma \right)\]
	is a continuous function of $(\omega, F, E_F)$ as claimed.
	
	Lastly, to check (c) consider any $\beta \in \Gamma$. Since left-multiplication by $\beta$ provides a bijection from the $E_F$-classes to the $E_{\beta^d \cdot F}$-classes, for each $C \in \Gamma / E_F$ we can use the point $\widehat \gamma_{\beta C} = \beta \gamma_C$ as the representative for $\beta C$ in computing $\theta(\beta^s \cdot \omega, \beta^d \cdot F)$. Then by (\ref{proofitem:fmapb}) and the fact that the map $f$ is $\Gamma^s$-equivariant we have
	\begin{align*}
		\beta^s_* \theta(\omega, F) & = \beta^s_* f_* \prod_{C \in \Gamma / E_F} \left( \prod_{\delta \in C} \sigma_F^\omega(\delta, \gamma_C) \right)_* \omega(\gamma_C)\\
		& = f_* \prod_{\beta C \in \Gamma / E_{\beta^d \cdot F}} \left( \prod_{\beta \delta \in \beta C} \sigma_F^\omega(\delta, \gamma_C) \right)_* \omega(\gamma_C)\\
		& = f_* \prod_{\beta C \in \Gamma / E_{\beta^d \cdot F}} \left( \prod_{\beta \delta \in \beta C} \sigma_{\beta^d \cdot F}^{\beta^s \cdot \omega}(\beta \delta, \beta \gamma_C) \right)_* \omega(\gamma_C)\\
		& = f_* \prod_{\beta C \in \Gamma / E_{\beta^d \cdot F}} \left( \prod_{\delta \in \beta C} \sigma_{\beta^d \cdot F}^{\beta^s \cdot \omega}(\delta, \hat \gamma_{\beta C}) \right)_* \omega(\beta^{-1} \hat \gamma_{\beta C})\\
		& = f_* \prod_{C \in \Gamma / E_{\beta^d \cdot F}} \left( \prod_{\delta \in C} \sigma_{\beta^d \cdot F}^{\beta^s \cdot \omega}(\delta, \hat \gamma_C) \right)_* (\beta^s \cdot \omega)(\hat \gamma_C) = \theta(\beta^s \cdot \omega, \beta^d \cdot F).\qedhere
	\end{align*}
\end{proof}

\end{document}
